
\DeclareMathOperator{\lcm}{lcm}

\newcommand{\pref}[1]{\operatorname{Pref}(#1)}
\newcommand{\suff}[1]{\operatorname{Suff}\left(#1\right)}
\newcommand{\factor}[1]{\operatorname{Fact}(#1)}
\newcommand{\fact}[1]{\ensuremath{\operatorname{Fact}(#1)}}
\newcommand{\constrsync}[1]{\text{\normalfont \textsc{Sync}\textsc{[}\mbox{$#1$}\textsc{]}}\xspace}
\newcommand{\concat}[2]{#2 #1}
\newcommand{\concatcirc}[2]{\ensuremath{#2 \circ #1}}
\newcommand{\applyfunc}[2]{#2(#1)}
\newcommand{\applytrans}[3]{\ensuremath{#1(#2,#3)}}
\newcommand{\applytransexplicit}[3]{\ensuremath{#1(#2,#3)\xspace\,}}
\newcommand{\semiautomat}[4]{\mathcal #1 = (#2, #3, #4)}
\newcommand{\automat}[6]{\ensuremath{\mathcal #1 = (#2, #3, #4, #5, #6)}}  % entfernt überstrich, wenn man es so benutzt: \automat{\overline{\mathcal A}_i}{Q_i \cup \{p_i\}}{\overline{\delta_i}}{p_i}{F_i}
\newcommand{\semiautomatexpl}[4]{\ensuremath{\mathcal #1 = (#2, #3, #4)}}
\newcommand{\automatp}[6]{\ensuremath{#1 = (#2, #3, #4, #5, #6)}}
\newcommand{\perm}{\operatorname{perm}}
\newcommand{\Syn}{\operatorname{Syn}}
\newcommand{\Cerny}{\v{C}ern\'{y}\xspace}
\newcommand{\orbit}[2]{\ensuremath{#2(#1)}}
\newcommand{\Orb}{\operatorname{Orb}}
\newcommand{\rt}[1]{\operatorname{rt}{(#1)}}

\newcommand{\rank}[1]{\ensuremath{\operatorname{rank}(#1)}}
\newcommand{\defect}[1]{\ensuremath{\operatorname{def}(#1)}} 
\newcommand{\Orbit}[2]{\ensuremath{\operatorname{Orb}_{#2}(#1)}}
\newcommand{\Stab}[2]{\ensuremath{{#2}_{#1}}}
\newcommand{\lereach}{\ensuremath{\precsim_R}}
\newcommand{\eqreach}{\ensuremath{\simeq_R}}
\newcommand{\lelogm}{\ensuremath{\le_m^{\log}}}
\newcommand{\lePm}{\ensuremath{\le_m^{\PTIME}}}
\newcommand{\leACZEROm}{\ensuremath{\le_m^{\AC^0}}}
\newcommand{\leACkm}[1]{\ensuremath{\le_m^{\AC^{#1}}}}
\newcommand{\leNCkm}[1]{\ensuremath{\le_m^{\NC^{#1}}}}
\newcommand{\leNCkT}[1]{\ensuremath{\le_T^{\NC^{#1}}}}
\newcommand{\eqlogm}{\ensuremath{\equiv_m^{\log}}}
\newcommand{\eqPm}{\ensuremath{\equiv_m^{\PTIME}}}
\newcommand{\eqACZEROm}{\ensuremath{\equiv_m^{\AC^0}}}
\newcommand{\eqACkm}[1]{\ensuremath{\equiv_m^{\AC^{#1}}}}
\newcommand{\eqNCkm}[1]{\ensuremath{\equiv_m^{\NC^{#1}}}}
\newcommand{\eqNCkT}[1]{\ensuremath{\equiv_T^{\NC^{#1}}}}
\newcommand{\fword}[2]{\ensuremath{{#1}_{#2}}}
\newcommand{\SI}[2]{\ensuremath{[#1 \mapsto #2]}}
\newcommand{\move}[1]{\operatorname{move}(#1)}
\newcommand{\fix}[1]{\operatorname{fix}(#1)}
\newcommand{\dupl}[1]{\operatorname{dupl}(#1)}
\newcommand{\excl}[1]{\operatorname{excl}(#1)}
\newcommand{\TM}[1]{\mathcal T_M(#1)}
\newcommand{\TS}[1]{\mathcal T_S(#1)}
\newcommand{\Z}{\mathbb{Z}}
\newcommand{\PA}[1]{\mathcal P(#1)}

\newcommand{\rYES}{{\sc YES}\xspace}
\newcommand{\rNO}{{\sc NO}\xspace}
\newcommand{\eqLp}[1]{\mathrel{\ensuremath{\equiv_{L}^{\ge #1}}}}
\newcommand{\eqLn}{\mathrel{\eqLp{n}}}
\newcommand{\eqLpp}[2]{\mathrel{\ensuremath{\equiv_{#1}^{\ge #2}}}}
\newcommand{\eqLL}[2]{\mathrel{\ensuremath{\equiv_{#1, #2}}}}
\newcommand{\NeqLL}[2]{\mathrel{\ensuremath{\not\equiv_{#1,#2}}}}
\newcommand{\eqLsimU}[3]{\mathrel{\ensuremath{\equiv_{#1}^{#2, #3}}}}
\newcommand{\preeqLL}[2]{\mathrel{\ensuremath{\sim_{#1, #2}}}}
 \newcommand{\excess}[2]{\operatorname{exc}_{#1}(#2)}
 
%\newcommand{\shift}[1]{\mathsf{X}_{#1}} % X_{#1}
%\newcommand{\term}[1]{\textit{#1}}
%\newcommand{\leftm}{{\vartriangleright}}  % "left marker"
%\newcommand{\rightm}{{\vartriangleleft}} % "right marker"

\newcommand{\Zb}{\mathbb{Z}}
\newcommand{\minaut}[1]{\ensuremath{\operatorname{MinAut}(#1)}}
\newcommand{\restrictAutGamS}[3]{\ensuremath{#1_{|#2, #3}}}
\newcommand{\upperset}[1]{\mbox{\ensuremath{\displaystyle  \uparrow \!\!#1}}}
\newcommand{\lowerset}[1]{\ensuremath{\downarrow \!\!#1}}
\renewcommand{\vec}[1]{\ensuremath{\boldsymbol{#1}}}


\newcommand{\PTIME}{\textsf{P}\xspace}
\newcommand{\LOGSPACE}{\textsf{L}\xspace}
\newcommand{\WNL}{\textsf{WNL}\xspace}
\newcommand{\WSYNC}{\textsf{W[Sync]}\xspace}

%\renewcommand{\SAT}{\textsc{\sc SAT}\xspace}
\newcommand{\SIS}{\textsc{\sc Sync\-Into}\xspace}
\newcommand{\ST}{\textsc{\sc Set\-Trans}\xspace}
\newcommand{\DST}{\textsc{\sc Disj\-Set\-Trans}\xspace}
\newcommand{\SFS}{\text{\sc Sync\-From}\xspace}
\newcommand{\DFAS}{\text{\sc Sync}\xspace}
\newcommand{\DFASWS}{\text{\sc Sync\-Short}\xspace}
\newcommand{\DFASS}{\text{\sc Sync\-State}\xspace}
\newcommand{\DFASW}{\text{\sc Sync\-Words}\xspace}
\newcommand{\DFAINT}{\text{\sc Intersection}\xspace}
\newcommand{\DAG}{\text{\sc DAG}\xspace} 
\newcommand{\DFAST}{\textsc{\sc Sync\-Trap}\xspace}
\newcommand{\THREESAT}{\text{\sc 3SAT}\xspace}
\newcommand{\ONEDAGST}{\text{\sc $1$-DAG-ST}\xspace}
\newcommand{\ONEDAG}{\text{\sc $1$-DAG}\xspace}
\newcommand{\TWODAG}{\text{\sc $2$-DAG}\xspace}
\newcommand{\TWODAGST}{\text{\sc $2$-DAG-ST}\xspace}
\newcommand{\KDAG}{\text{\sc $k$-DAG}\xspace}
\newcommand{\SCP}{\text{\sc Single\-Cycle}\xspace}
\newcommand{\PDFAINT}{\text{\sc PDFA-Intersection}\xspace}
\newcommand{\SCC}{\text{\sc Strongly\-Conn}\xspace}
\newcommand{\kDAGST}[1]{\text{\sc $#1$-DAG-ST}\xspace}
\newcommand{\KDAGST}{\kDAGST{k}\xspace}
\newcommand{\ORD}{{\sc OrderBetweenVertices}\xspace}
\newcommand{\REACH}{\text{\sc Reach}\xspace}
\newcommand{\REACHF}{\text{\sc ReachFrom}\xspace}

\newcommand{\problemdef}[4]{
	% \problemdef{title}{non-small-caps reference/comment/... in title}{input}{problem}
	% Next line: width of the drawn box, spacing left (right with "right = ...")
	\begin{tcolorbox}[width = \textwidth,colback=white,arc=0pt,outer arc=0pt,boxrule=0.7pt,left =0.5em,right=0em]
		\textsc{#1} #2
		\\[2pt]
		% Next line: table width
		\begin{tabular}{ @{}l p{0.84\textwidth} c }
			\textsf{Input:} & #3 \\[.5pt]
			\textsf{Problem:} & #4
		\end{tabular}
		\vspace{-0.25em}
	\end{tcolorbox}
}

\crefname{lemma}{Lemma}{Lemmas}
\crefname{proposition}{Proposition}{Propositions}
\crefname{theorem}{Theorem}{Theorems}
\crefname{equation}{Equation}{Equations}
\crefname{item}{Item}{Items}
\crefname{corollary}{Corollary}{Corollaries}
\crefname{section}{Section}{Sections}
\crefname{example}{Example}{Examples}
\crefname{Fact}{Fact}{Facts}
\crefname{table}{Table}{Tables}
\crefname{remark}{Remark}{Remarks}
\crefname{part}{Part}{Parts}
\crefname{figure}{Figure}{Figures}
\crefname{chapter}{Chapter}{Chapters}
\crefname{definition}{Definition}{Definitions}
\crefname{corollary}{Corollary}{Corollaries}

% https://tex.stackexchange.com/questions/410585/cleveref-format-for-enumerate
\crefname{enumi}{Item}{Items}
\crefname{enumii}{Item}{Items}
\crefname{enumiii}{Item}{Items}

% https://latex.org/forum/viewtopic.php?t=6667
\newcommand{\crefitem}[2]{\cref{#1},~Item~\ref{#2}}
\newcommand{\Crefitem}[2]{\Cref{#1},~Item~\ref{#2}}


%%\newcommand{\THM}{Thm.}
%%\newcommand{\COR}{Cor.}
%%\newcommand{\SEC}{Sect.}
%%\newcommand{\PROP}{Prop.}
%%\newcommand{\REM}{Rem.}
%%\newcommand{\LEM}{Lem.}
%%\newcommand{\EQN}{Eqn.}
%%\newcommand{\FIG}{Fig.}
\newcommand{\THM}{Theorem}
\newcommand{\COR}{Corollary}
\newcommand{\PROP}{Proposition}
%\newcommand{\REM}{Remark}
\newcommand{\SEC}{Section}
\newcommand{\LEM}{Lemma}
%\newcommand{\EQN}{Equation}
\newcommand{\FIG}{Figure}
%\newcommand{\EX}{Example}
%\newcommand{\CHAP}{Chapter}





%% https://tex.stackexchange.com/questions/570985/how-to-use-a-fact-as-an-equation-in-latex
%\newcommand\myparbox[1]{\parbox{0.667\textwidth}{\raggedright #1}} % choose width suitably
%
%
%
%
%% \sqcup  oder \square
%% https://tex.stackexchange.com/questions/136485/empty-clause-symbol
%\newcommand{\TMblank}{%
%\text{\fboxsep=-.2pt\fbox{\rule{0pt}{1.35ex}\rule{1.15ex}{0pt}}}%
%}

\tikzstyle{every state}=[blur shadow,fill={rgb:black,0.4;white,10}]


% Nur fuer pdflatex
%\DeclareUnicodeCharacter{0301}{*************************************}
\DeclareUnicodeCharacter{0301}{\i}








\DeclareSymbolFont{rsfscript}{OMS}{rsfs}{m}{n}
\DeclareSymbolFontAlphabet{\mathrsfs}{rsfscript}

% https://tex.stackexchange.com/questions/22012/hyperref-and-imakeidx-do-not-work-together
\newcommand{\BHT}[1]{\normalsize{\hyperpage{#1}}\normalsize}
\newcommand{\IN}[1]{\index{#1|BHT}}
%\renewcommand{\index[1]{\index{#1|BH}}

% Mit \IN geht das nicht
%\IN{function composition|see{composition}}
%\index{function composition|see{composition}}






% https://tex.stackexchange.com/questions/21234/doing-something-only-when-the-draft-option-is-on
\makeatletter
\def\ifdraft{\ifdim\overfullrule>\z@
	\expandafter\@firstoftwo\else\expandafter\@secondoftwo\fi}
\makeatother

% https://tex.stackexchange.com/questions/32208/footnote-runs-onto-second-page
\interfootnotelinepenalty=10000 %% Completely prevent breaking of footnotes





% https://tex.stackexchange.com/questions/655563/what-is-the-meaning-of-no-math-alphabet-change-to-frozen-version-normal-on-inpu
\def\@font@info#1{}
\setcounter{localmathalphabets}{0}





% Andrew Stacey's code from
% https://tex.stackexchange.com/a/50054/3954
%\makeatletter
%\tikzset{%
%	remember picture with id/.style={%
%		remember picture,
%		overlay,
%		save picture id=#1,
%	},
%	save picture id/.code={%
%		\edef\pgf@temp{#1}%
%		\immediate\write\pgfutil@auxout{%
%			\noexpand\savepointas{\pgf@temp}{\pgfpictureid}}%
%	},
%	if picture id/.code args={#1#2#3}{%
%		\@ifundefined{save@pt@#1}{%
%			\pgfkeysalso{#3}%
%		}{
%			\pgfkeysalso{#2}%
%		}
%	}
%}

%\def\savepointas#1#2{%
%	\expandafter\gdef\csname save@pt@#1\endcsname{#2}%
%}
%
%\def\tmk@labeldef#1,#2\@nil{%
%	\def\tmk@label{#1}%
%	\def\tmk@def{#2}%
%}
%
%\tikzdeclarecoordinatesystem{pic}{%
%	\pgfutil@in@,{#1}%
%	\ifpgfutil@in@%
%	\tmk@labeldef#1\@nil
%	\else
%	\tmk@labeldef#1,(0pt,0pt)\@nil
%	\fi
%	\@ifundefined{save@pt@\tmk@label}{%
%		\tikz@scan@one@point\pgfutil@firstofone\tmk@def
%	}{%
%		\pgfsys@getposition{\csname save@pt@\tmk@label\endcsname}\save@orig@pic%
%		\pgfsys@getposition{\pgfpictureid}\save@this@pic%
%		\pgf@process{\pgfpointorigin\save@this@pic}%
%		\pgf@xa=\pgf@x
%		\pgf@ya=\pgf@y
%		\pgf@process{\pgfpointorigin\save@orig@pic}%
%		\advance\pgf@x by -\pgf@xa
%		\advance\pgf@y by -\pgf@ya
%	}%
%}
%\newcommand\tikzmark[2][]{%
%	\tikz[remember picture with id=#2] {#1;}}
%\makeatother
%% end of Andrew's code
%
%\newcommand\ShadeCell[4][0pt]{%
%	\begin{tikzpicture}[overlay,remember picture]%
%		\shade[#4] ( $ (pic cs:#2) + (0pt,1.9ex) $ ) rectangle ( $ (pic cs:#3) + (0pt,-#1*\baselineskip-.8ex) $ );
%	\end{tikzpicture}%
%}%
%
%\newcommand\HatchedCell[4][0pt]{%
%	\begin{tikzpicture}[overlay,remember picture]%
%		\fill[#4] ( $ (pic cs:#2) + (0,1.9ex) $ ) rectangle ( $ (pic cs:#3) + (0pt,-#1*\baselineskip-.8ex) $ );
%	\end{tikzpicture}%
%}%





% https://tex.stackexchange.com/questions/161338/can-cleveref-be-made-to-use-the-oxford-comma-for-multiple-citations
\newcommand{\creflastconjunction}{, and\nobreakspace}


%		\tikzset{every state/.style={minimum size=1pt}}


\newcommand{\tikzcell}[1]{%
	\begin{tikzpicture}[
		baseline=(current bounding box.center),
		every node/.style={align=center},
		inner sep=0pt,
		outer sep=0pt,
		overlap/.style={draw, minimum height=6mm, minimum width=1.6cm},
		normal/.style={draw, minimum height=6mm, minimum width=1.6cm}
		]
		#1
	\end{tikzpicture}%
}


% Mark thesis as confidential
% \usepackage[angle=0,fontsize=16pt,color=black,pos={20.5cm, 0.5cm} ,anchor=rt,alignment=r,text=TAJEMNICA\ PŁ\\Egz.\ nr\ X]{draftwatermark}

% Highlight thesis author in bibliobraphy
% \usepackage{xstring}
% \let\originalbibitem\bibitem
% \def\bibitem#1#2\par{%
	%   \noexpandarg
	%   \originalbibitem{#1}
	%   \StrSubstitute{#2}{I.Nazwisko}{\textbf{I.Nazwisko}}\par}





\newcommand{\cg}{\cellcolor{gray!50}}
%\newcommand{\cd}{\cellcolor{gray!20}}
\newcommand{\cd}{}


\def\ygap{2}
\def\xgap{2}
\def\arrowgap{.125}
\tikzset{
	arc center/.is choice,
	arc center/left/.code={\def\isLeft{1}\tikzset{arc to}},
	arc center/right/.code={\def\isLeft{0}\tikzset{arc to}},
	arc center offset/.code={\def\centerOffset{#1}\tikzset{arc to}},
	arc direction/.is choice,
	arc direction/cw/.code={\def\isCw{1}\tikzset{arc to}},
	arc direction/ccw/.code={\def\isCw{0}\tikzset{arc to}},
	arc to/.style={
		to path={
			[evaluate={
				coordinate \u, \v, \m, \mv, \mvOrtho, \c, \cu, \cv;
				\u = (\tikztostart);
				\v = (\tikztotarget);
				\m = (\u)!0.5!(\v);
				\mv = (\v)-(\m);
				if \isLeft then {
					\mvOrtho = (-\mvy, \mvx);
				} else {
					\mvOrtho = (\mvy, -\mvx);
				};
				\c = (\m)+(\mvOrtho);
				\c = (\m)!\centerOffset!(\c);
				\cu = (\u)-(\c);
				\cv = (\v)-(\c);
				\radius = veclen(\cu);
				\startAngle = atan2(\cuy, \cux);
				\endAngle = atan2(\cvy, \cvx);
				if \isCw then {
					if \startAngle < \endAngle then {
						\endAngle = \endAngle - 360;
					};
				} else {
					if \startAngle > \endAngle then {
						\endAngle = \endAngle + 360;
					};
				};
			}]
			arc[radius=\radius pt, start angle=\startAngle, end angle=\endAngle] \tikztonodes
		}
	},
	vertex/.style={circle, inner sep=1pt},
	left loop/.style={arc center=left, arc direction=cw, arc center offset=10pt},
	right loop/.style={arc center=right, arc direction=ccw, arc center offset=10pt},
	invisible/.style={text opacity=0},
	-{Stealth[round, flex=1]}
}








\newcommand{\decisionproblem}[3]{%
	\par\addvspace{.5\baselineskip}
	\noindent
	
	\begin{tabularx}{\textwidth}{@{\hspace{\parindent}} l X c}
		\multicolumn{2}{@{\hspace{\parindent}}l}{#1} \\
		\textsf{Input:}    & #2 \\
		\textsf{Question:} & #3
	\end{tabularx}
	
	\par\addvspace{.5\baselineskip}
}



\NewDocumentCommand{\VertexChain}{O{} m m m m m}{%
	% #1 = optionaler Startknoten
	% #2 = prefix
	% #3,#4,#5,#6 = Indizes
	
	% --- Start ---
	\IfBlankTF{#1}{%
		\node[state,label=above:$v_{#3}$] (#2-#3) {};
	}{%
		\node[state,label=above:$v_{#3}$, right=of #1, xshift=0.5em] (#2-#3) {};
	}
	
	\node[right=of #2-#3] (#2-d1) {$\cdots$};
	
	\node[state,
	label={[xshift=-7pt]above:$v_{#4}$}, 
	right=of #2-d1
	] (#2-#4) {};
	
	\node[state, label=above:$v_{{#4}+1}$, right=of #2-#4] (#2-#4n) {};	
	
	\node[right=of #2-#4n] (#2-d2) {$\cdots$};
	
	\node[state,
	label={[xshift=-7pt]above:$v_{#5}$}, 
	right=of #2-d2] (#2-#5) {};
	
	\node[state,label=above:$v_{{#5}+1}$, right=of #2-#5] (#2-#5n) {};
	
	\node[right=of #2-#5n] (#2-d3) {$\cdots$};
	
	\node[state,
	label={[xshift=-7pt]above:$v_{#6}$},
	right=of #2-d3
	] (#2-#6) {};
	
	\path[->]
	(#2-#3) edge (#2-d1)
	(#2-d1) edge (#2-#4)
	%(#2-#4) edge (#2-#4n)
	(#2-#4n) edge (#2-d2)
	(#2-d2) edge (#2-#5)
	%(#2-#5) edge (#2-#5n)
	(#2-#5n) edge (#2-d3)
	(#2-d3) edge (#2-#6);
}

\newcommand{\thesistitle}{The Set of Synchronizing Words and the Constrained Synchronization Problem}

% https://tex.stackexchange.com/questions/20576/passing-newenvironment-parameters-to-the-end-block
\NewDocumentEnvironment{proofitem}{m}
{%
	\noindent\emph{Proof of #1:}%
}
{%
	\emph{[End, Proof of #1]}%
}